\section{Design of a dynamic current allocator for the TCV shape controller}
\label{sec:tcvalldesign}
The design of a dynamic coil current allocator for the TCV shape controller requires an inspection of the preexisting control scheme, to clarify which among the novel formulations presented in \Cref{sec:all_design} is the most appropriate. Subsequently, the selected solution is adapted to the specific case of the TCV shape controller, accounting for any differences arising with respect to the previously analyzed scenarios.

\subsection{Problem statement}
\label{sec:tcvalldesign_problem}
Consider the simplified control scheme given in \Cref{fig:tcvall}, which summarizes the interconnection of the main control subsystems with the TCV plant. Before going further, it is instrumental to anticipate that every dynamic model considered in the following is continuous-time, linear time-invariant; this is possible thanks to the powerful yet real-time equilibrium reconstruction codes that are part of the SCD control suite, which allow control engineers to work on a linearized model of the plant around the prescribed shape equilibrium (see \Cref{sec:tcv,sec:scd,sec:sota_tcvshapectrl}).
\begin{figure}[t]
  \centering
  \includesvg[width=\linewidth]{images/ch7/tcvall.svg}
  \caption{Simplified diagram of the TCV control scheme: the $P_h$ block represents the actual plant, couple with the \emph{hybrid} feedback controller $C_h$; the shape controller is represented by the $C_sh$ block, whereas the proposed allocator is embedded in the $\mathcal{A}\ell\ell$ block.}
  \label{fig:tcvall}
\end{figure}

The central block $P_h$ represents the dynamics of the plasma and coils, described by the following equations:
\begin{subequations}
  \label{eq:tcv_ph}
  \begin{align}
    \dot{x}_{ph}(t) &= A_{ph}x_{ph}(t) + B_{ph}u_{h}(t), & t\in\R_{\geq 0},\\
    y(t) &= C_{ph}x_{ph}(t) + D_{ph}u_{h}(t), & t\in\R_{\geq 0},
  \end{align}
\end{subequations}
where $x_{ph}(t)\in\R^{275}$ is the state vector at time $t$, containing the currents in the active circuits $I_a(t)\in\R^{19}$ and in the filaments used to represent the vessel, namely $I_u(t)\in\R^{256}$; $u_{h}(t)\in\R^{19}$ is the plant input, containing the voltages to be applied to the active circuits $V_a(t)\in\R^{19}$; $y(t) = \begin{bmatrix}y_h(t)' & y_{sh}(t)' & I_{EF}(t)'\end{bmatrix}'$ is the plant output, containing the controlled variables $y_h(t)\in\R^{24}$, namely the plasma current and its radial and vertical positions as well as the currents of the PF coils, the shape output $y_{sh}(t)\in\R^{\bboutodim}$ whose dimension depends on the design of the shape controller for the desired equilibrium, and the measurements of the currents of the EF coils, namely $I_{EF}(t)\in\R^{16}$.

The inner control loop comprises also the feedback controller $C_h$, informally termed \emph{hybrid} due to its composition of both analog and digital devices \cite{lister_control_1997}. The role of $C_h$ is to compute the active voltages $V_a$ which are to be applied to the PF coils to track the reference input $r_h$, taking care of the vertical and radial stabilization of the plasma and controlling the current flow in the PF coils. Its state-space representation is given by
\begin{subequations}
  \label{eq:tvc_ch}
  \begin{align}
    \dot{x}_{ch}(t) &= A_{ch}x_{ch}(t) + B_{ch}e_h(t), & t\in\R_{\geq 0},\\
    u_h(t) &= C_{ch}x_{ch}(t) + D_{ch}e_h(t), & t\in\R_{\geq 0},
  \end{align}
\end{subequations}
where $x_{ch}(t)\in\R^{n_{ch}}$ is the state vector of the hybrid controller at time $t$, $e_h(t)\in\R^{24}$ is the hybrid error, $e_h = r_h - y_h$, and $u_h(t)\in\R^{19}$ is the hybrid controller output.

The shape controller is embedded in the $C_{sh}$ block, and it is given by the following state-space representation:
\begin{subequations}
  \label{eq:tcv_csh}
  \begin{align}
    \dot{x}_{csh}(t) &= A_{csh}x_{csh}(t) + B_{csh}e_{sh}(t), & t\in\R_{\geq 0},\\
    u_{sh}(t) &= C_{csh}x_{csh}(t) + D_{csh}e_{sh}(t), & t\in\R_{\geq 0},
  \end{align}
\end{subequations}
where $x_{csh}(t)\in\R^{n_{csh}}$ is the state vector of the shape controller at time $t$, $e_{sh}(t)\in\R^{\bboutodim}$ is the shape error, $e_{sh} = r_{sh} - y_{sh}$, and $u_{sh}(t)\in\R^{24}$ is the shape controller output, namely the variations to be applied to modify the prescribed feedforward references $r_{h,ff}$ for the hybrid controller, which are precomputed for each shot.

\begin{assumption}
  \label{ass:tcvstab}
  The closed-loop system \eqref{eq:tcv_ph}--\eqref{eq:tcv_csh}, interconnected as in
  \begin{align}
    \label{eq:tcv_interconn}
    e_{sh} &= r_{sh} - y_{sh}, & r_h &= r_{h,ff} + u_{sh},
  \end{align}
  is well-posed and aymptotically stable.\closestatement
\end{assumption}

% ass: overact
% ass: constant shape references
% rem: hybrid ff given, not constant in practice, some ramp, cannot really circumvent all the time
% define steady state cost function
% prb: steady state cost, full output invisibility
